% ============================================================
% Conteúdo acadêmico puro do documento "Requisitos Funcionais e
% Não Funcionais". SEM \documentclass, \begin{document}/\end{document}
% nem título — pensado para ser reaproveitado via \input{} tanto no
% documento modular quanto, futuramente, no TCC principal.
% ============================================================

\section{Apresentação}

O RadarTorres é um protótipo que detecta alvos por sensoriamento, rastreia-os em um radar
circular, seleciona automaticamente a torre mais adequada e permite o acionamento manual ou
automático de um indicador demonstrativo. Este documento reúne os requisitos funcionais (o que o
sistema deve fazer, do ponto de vista de um ator) e os requisitos não funcionais (que
características de qualidade essa funcionalidade deve ter) levantados a partir do código-fonte
(\texttt{src/RadarTorres.App/}), do firmware (\texttt{Arduino/ArduinoSimulation.ino}) e da
documentação técnica já existente no projeto.

\textbf{Convenções:} a rastreabilidade completa de cada item até classe/método do código, bem
como o status de implementação de cada requisito, está em \texttt{Matriz\_de\_Rastreabilidade.md}
(nesta mesma pasta), para não duplicar aqui.

\section{Requisitos Funcionais}

\begin{longtable}{p{1.3cm}p{8.0cm}p{3.6cm}}
  \caption{Requisitos funcionais do RadarTorres} \label{tab:rf} \\
  \hline
  \textbf{ID} & \textbf{Requisito} & \textbf{Ator(es)} \\
  \hline
  \endfirsthead
  \multicolumn{3}{c}{\small \tablename~\thetable\ (continuação)} \\
  \hline
  \textbf{ID} & \textbf{Requisito} & \textbf{Ator(es)} \\
  \hline
  \endhead
  \hline
  \endfoot
  \small
  RF01 & Comunicação serial com o Arduino & Usuário, Arduino \\
  RF02 & Recepção e validação de leituras de alvo & Arduino, Sistema \\
  RF03 & Rastreamento de alvos em tempo real & Sistema \\
  RF04 & Exibição do radar circular em tempo real & Usuário \\
  RF05 & Seleção automática de torre & Sistema \\
  RF06 & Acionamento demonstrativo automático (D1) & Sistema \\
  RF07 & Modos de operação (Verde/Amarelo/Vermelho) (D1) & Usuário \\
  RF08 & Autenticação multiusuário & Usuário \\
  RF09 & Controle de acesso por perfil & Sistema, Usuário \\
  RF10 & Troca de senha & Usuário \\
  RF11 & Registro histórico de detecções & Sistema \\
  RF12 & Visualização de objetos detectados & Usuário \\
  RF13 & Exportação de objetos detectados (CSV/XML/PDF) & Usuário \\
  RF14 & Importação de objetos detectados (CSV/XML) & Usuário \\
  RF15 & Auditoria de ações realizadas & Sistema, Usuário \\
  RF16 & Auditoria de alterações de modo & Sistema, Usuário \\
  RF17 & Gerenciamento de usuários & Usuário (Admin) \\
  RF18 & Preferências de usuário (idioma, tema, sidebar) & Usuário \\
  RF19 & Abertura de chamado de ajuda & Usuário \\
  RF20 & Tratamento administrativo de chamados & Usuário (Admin) \\
  RF21 & Painel principal com cards personalizáveis & Usuário \\
  RF22 & Console de eventos fixável na lateral & Usuário \\
  RF23 & Gestão de zonas mortas (áreas de exclusão) & Usuário (Admin), Sistema \\
  RF24 & Detecção e compilação de sketch Arduino pela interface & Usuário, Sistema \\
  RF25 & Monitor serial e persistência de configurações da aba Arduino & Usuário, Sistema, Arduino \\
\end{longtable}

\textbf{D1} = divergência entre especificação e implementação atual (o código ainda expõe um
caminho de acionamento manual e uma nomenclatura de modos diferente da especificada) — detalhada
em \texttt{Limitacoes\_Conhecidas.md}.

\section{Requisitos Não Funcionais}

\begin{longtable}{p{1.6cm}p{8.4cm}p{3.4cm}}
  \caption{Requisitos não funcionais do RadarTorres} \label{tab:rnf} \\
  \hline
  \textbf{ID} & \textbf{Requisito} & \textbf{Categoria} \\
  \hline
  \endfirsthead
  \multicolumn{3}{c}{\small \tablename~\thetable\ (continuação)} \\
  \hline
  \textbf{ID} & \textbf{Requisito} & \textbf{Categoria} \\
  \hline
  \endhead
  \hline
  \endfoot
  \small
  RNF01 & Redesenho do radar em $\sim$150\,ms & Desempenho \\
  RNF02 & Leitura serial não bloqueante & Desempenho \\
  RNF03 & Degradação graciosa de falhas de hardware/serial & Confiabilidade \\
  RNF04 & Distância mínima de segurança obrigatória no acionamento & Segurança \\
  RNF05 & Acionamento é só demonstrativo, nunca armamento real & Segurança \\
  RNF06 & Senha de usuário nunca em texto puro & Segurança \\
  RNF07 & Zonas mortas só editáveis por Administrador & Segurança \\
  RNF08 & Montagem segura de comandos de processo externo & Segurança \\
  RNF09 & Arduino CLI nunca baixado automaticamente & Segurança \\
  RNF10 & Interface responsiva a mudanças de resolução/tamanho & Usabilidade \\
  RNF11 & Troca de idioma sem reinício & Usabilidade \\
  RNF12 & Consistência visual entre temas & Usabilidade \\
  RNF13 & Limite de linhas nos consoles (anti-crescimento de memória) & Desempenho \\
  RNF14 & Compilação do sketch não bloqueia a interface & Desempenho \\
  RNF15 & Runtime portátil, sem dependência externa & Portabilidade \\
  RNF16 & Upgrade preserva configurações do usuário & Manutenibilidade \\
  RNF17 & Persistência gravável sem privilégio de administrador & Manutenibilidade \\
  RNF18 & Cobertura de testes automatizados dos componentes críticos & Confiabilidade \\
  RNF19 & Plataforma-alvo Windows Desktop & Compatibilidade \\
  RNF20 & Tolerância a mensagens malformadas do Arduino & Confiabilidade \\
\end{longtable}
